% -------------------------------------------------------------------------
% WBS ID: RES-NUM-TIM
% Deliverable: Time Integration, CFL Control and Stiffness Treatment
% Status: Draft complete
% Evidence status: Regional and local timestep constraints established at research level
% Review status: Not reviewed
% Last updated: 2026-07-18
% Dependencies: RES-NUM-FRM, RES-NUM-FLX, RES-NUM-SRC, RES-NUM-WET, RES-NUM-SPEC
% Downstream interfaces: RES-VER-CNV, Software time-integration implementation
% -------------------------------------------------------------------------

\section{RES-NUM-TIM --- Time Integration, CFL Control and Stiffness Treatment}
\label{sec:res-num-tim}

\subsection{Purpose and Scope}
\label{sec:res-num-tim-purpose}

% TODO: Define what this Deliverable covers and explicitly excludes.

\subsection{Research Questions}
\label{sec:res-num-tim-research-questions}

% TODO: State the questions that must be answered before the Deliverable
% can be accepted.

\subsection{Terminology and Definitions}
\label{sec:res-num-tim-definitions}

% TODO: Define the technical terminology, symbols, quantities and
% classifications used in this Deliverable.

\subsection{Literature Evidence}
\label{sec:res-num-tim-literature-evidence}

% TODO: Record evidence from peer-reviewed papers, authoritative datasets,
% standards, technical reports and primary documentation.
%
% Do not create a detached literature-summary list. Explain how each source
% supports, contradicts or limits a project decision.

\subsection{Governing Relations and Quantitative Definitions}
\label{sec:res-num-tim-governing-relations}

% TODO: Record relevant equations, dimensionless groups, empirical models,
% extraction rules or quantitative definitions.
%
% Use align environments for displayed equations.
% Every displayed equation must have a unique \label{eq:...}.
% Do not place punctuation after displayed equations.

\subsection{Discrete Formulation}
\label{sec:res-num-tim-discrete-formulation}

The regional time update is SSPRK(3,3) for the transport/source
residual excluding Manning friction, combined with Strang-split
semi-implicit Manning friction:
\cref{eq:res-num-spec-regional-strang-ssprk3}. Full IMEX RK is not
selected for the regional baseline because the retained stiff term is
cell-local Manning friction.

The local URANS--VOF model uses a segregated transient update. Within
each accepted time step the VOF transport, mixture-property update,
momentum predictor, pressure correction, turbulence equations and
metric extraction are ordered as specified in `RES-NUM-SPEC`.

\subsection{Conservation Properties}
\label{sec:res-num-tim-conservation-properties}

% TODO: Complete this RES-NUM Domain-specific subsection.

\subsection{Consistency and Formal Accuracy}
\label{sec:res-num-tim-consistency-formal-accuracy}

% TODO: Complete this RES-NUM Domain-specific subsection.

\subsection{Stability Requirements}
\label{sec:res-num-tim-stability-requirements}

% TODO: Complete this RES-NUM Domain-specific subsection.

\subsection{Mesh and Time-Step Requirements}
\label{sec:res-num-tim-mesh-time-step-requirements}

The regional timestep is controlled by the face signal speed and cell
geometry in
\cref{eq:res-num-spec-regional-face-signal-speed,eq:res-num-spec-regional-cell-timestep,eq:res-num-spec-regional-global-timestep}.
The provisional CFL value is $C_{\mathrm{CFL}}=0.5$, subject to
positivity and wet/dry verification.

The local timestep shall be the minimum of flow-advection, VOF
interface, gravity-wave and diffusion constraints:

\begin{align}
  \Delta t_{\mathrm{local}}
  &=
  \min_i
  \left(
    \Delta t_{\mathrm{CFL},i},\,
    \Delta t_{\alpha,i},\,
    \Delta t_{\mathrm{g},i},\,
    \Delta t_{\nu,i}
  \right)
  \label{eq:res-num-tim-local-timestep-minimum}
\end{align}

The local flow Courant constraint is:

\begin{align}
  \Delta t_{\mathrm{CFL},i}
  &=
  C_{\mathrm{CFL}}
  \frac{
    V_i
  }{
    \sum_{f\in\partial K_i}
    \left|
      \phi_f
    \right|
    +
    \epsilon_{\phi}
  }
  \label{eq:res-num-tim-local-flow-cfl}
\end{align}

The VOF/interface Courant constraint is:

\begin{align}
  \Delta t_{\alpha,i}
  &=
  C_{\alpha}
  \frac{
    V_i
  }{
    \sum_{f\in\partial K_i}
    \left|
      \phi_{\alpha,f}
    \right|
    +
    \epsilon_{\phi}
  }
  \label{eq:res-num-tim-local-interface-cfl}
\end{align}

The local gravity-wave constraint is:

\begin{align}
  \Delta t_{\mathrm{g},i}
  &=
  C_{\mathrm{g}}
  \frac{
    \Delta_i
  }{
    \sqrt{
      g h_i
    }
    +
    \epsilon_{\mathrm{g}}
  }
  \label{eq:res-num-tim-local-gravity-wave-cfl}
\end{align}

The diffusion constraint is:

\begin{align}
  \Delta t_{\nu,i}
  &=
  C_{\nu}
  \frac{
    \Delta_i^2
  }{
    \nu_{\mathrm{eff},i}
    +
    \epsilon_{\nu}
  }
  \label{eq:res-num-tim-local-diffusion-constraint}
\end{align}

where $\Delta_i$ is a representative local cell length and
$\nu_{\mathrm{eff},i}=\mu_{\mathrm{eff},i}/\rho_i$. The constants
$C_{\mathrm{CFL}}$, $C_{\alpha}$, $C_{\mathrm{g}}$ and $C_{\nu}$ shall
be verified for the selected VOF scheme, pressure-correction sequence,
mesh and wall treatment.

\subsection{Numerical Failure Modes}
\label{sec:res-num-tim-numerical-failure-modes}

% TODO: Complete this RES-NUM Domain-specific subsection.

\subsection{Verification Requirements}
\label{sec:res-num-tim-verification-requirements}

% TODO: Complete this RES-NUM Domain-specific subsection.

\subsection{Candidate Approaches}
\label{sec:res-num-tim-candidate-approaches}

% TODO: Define the credible candidate approaches considered.

\subsection{Critical Comparison}
\label{sec:res-num-tim-critical-comparison}

% TODO: Compare candidates using physically and project-relevant criteria.
% Do not select an approach solely on popularity or implementation ease.

\subsection{Assumptions and Applicability}
\label{sec:res-num-tim-assumptions}

% TODO: State assumptions and the conditions under which they are valid.

\subsection{Limitations and Failure Conditions}
\label{sec:res-num-tim-limitations}

% TODO: State where the evidence, model or selected approach ceases to be
% reliable or applicable.

\subsection{Project-Specific Conclusions}
\label{sec:res-num-tim-conclusions}

% TODO: Convert the evidence into conclusions specific to the Studentship
% project. Do not merely restate literature findings.

\subsection{Selected Approach and Rationale}
\label{sec:res-num-tim-selected-approach}

% TODO: Record the selected approach only after sufficient evidence exists.
% State the alternatives rejected and the reasons for rejection.

The selected regional time-integration approach is SSPRK(3,3) transport
with Strang-split Manning friction. This preserves an explicit
shock-capturing transport update while treating the friction damping in
a stable cell-local substep.

The selected local time-control approach is a segregated transient
URANS--VOF update with timestep selection governed by
\cref{eq:res-num-tim-local-timestep-minimum}. The final timestep is an
accepted-time-step value, not merely a requested value; failed
boundedness or convergence checks may force timestep rejection.

\subsection{Unresolved Decisions}
\label{sec:res-num-tim-unresolved-decisions}

% TODO: Record unresolved questions, required evidence and decision owners.

The timestep method is selected. The final CFL constant, local
subcycling policy, checkpoint interval and failure-recovery strategy are
implementation/validation decisions.

For the local model, the open timestep decisions are final Courant
limits, pressure-corrector count, turbulence subcycling if required,
adaptive timestep-reduction factor and coupling-history sampling
interval.

\subsection{Requirements and Handoffs}
\label{sec:res-num-tim-requirements}

% TODO: State requirements passed to other Research Domains, Software,
% verification activities or later execution work.

\subsection{Deliverable Completion Record}
\label{sec:res-num-tim-completion-record}

% TODO: Record whether the Deliverable is:
%
% - Not started
% - In progress
% - Draft complete
% - Reviewed
% - Accepted
%
% Acceptance requires:
%
% - the research questions to be answered;
% - claims to be supported by appropriate evidence;
% - assumptions and limitations to be explicit;
% - the project-specific conclusion to be stated;
% - downstream requirements to be traceable;
% - unresolved decisions to be closed or formally deferred.
